% BHU PYQ -- Transcribed & Corrected
% Paper: VDVMD-31 / VDVMD31: Vedanta
% Exam: Shastri (Hons.) / B.A. (Hons.) Semester III Examination, 2025-26 (NEP)
% Category: Multidisciplinary Course (MDC)
% Department: Department of Vedanta, Faculty of SVDV / Faculty of Arts, Banaras Hindu University

\documentclass[11pt,a4paper]{article}
\usepackage[a4paper,top=2.5cm,bottom=2.5cm,left=2.8cm,right=2.8cm,headheight=14pt]{geometry}
\usepackage{amsmath,amssymb}
\usepackage[T1]{fontenc}
\usepackage[utf8]{inputenc}
\usepackage{lmodern,microtype,fancyhdr,enumitem,hyperref,lastpage,parskip}

\hypersetup{
    colorlinks=true,
    urlcolor=black,
    linkcolor=black,
    pdftitle={VDVMD31: Vedanta -- BHU Shastri / B.A. Sem III 2025-26 (NEP)}
}

\pagestyle{fancy}
\fancyhf{}
\renewcommand{\headrulewidth}{0.4pt}
\renewcommand{\footrulewidth}{0.4pt}
\fancyhead[L]{\small \textit{Banaras Hindu University}}
\fancyhead[R]{\small \textit{VDVMD-31: Vedanta (Sem III 2025-26)}}
\fancyfoot[C]{\small Page \thepage\ of \pageref{LastPage}}

\newlist{parts}{enumerate}{2}
\setlist[parts,1]{label=(\alph*),leftmargin=*,itemsep=4pt,topsep=2pt}

\newcommand{\pts}[1]{\hfill \textit{[#1]}}

\begin{document}

\begin{center}
  {\large\bfseries BANARAS HINDU UNIVERSITY}\\[2pt]
  {\normalsize Shastri (Hons.) / B.A. Semester III Examination, 2025-26 (NEP)}\\[6pt]
  \rule{0.8\linewidth}{0.5pt}\\[6pt]
  {\large\bfseries DEPARTMENT OF VEDANTA}\\[4pt]
  {\large\bfseries Paper Code: VDVMD-31 : वेदान्त (Vedanta)}\\[2pt]
  {\small Ref. Code: 82/2011-III/59 / 2025-26/D-III/059}\\[4pt]
  {\normalsize Time: 2:30 Hours \hfill Full Marks: 70}
\end{center}

\rule{\linewidth}{0.4pt}

\smallskip

\noindent \textit{(Write your Roll No. at the top immediately on the receipt of this question paper)}\\[2pt]
\noindent \textit{सूचना : सभी प्रश्नों के उत्तर निर्देशानुसार दीजिए।}

\bigskip

\begin{center}
  {\large\bfseries SECTION A / खण्ड-अ}\\[2pt]
  {\normalsize (Long Answer Type Questions / दीर्घ उत्तरीय प्रश्न)}
\end{center}

\noindent \textit{सूचना : किन्हीं तीन प्रश्नों के उत्तर दीजिए। प्रत्येक प्रश्न 10 अंकों का है।}\\[2pt]

\bigskip

\noindent \textbf{Question 1.} वेदान्त दर्शन के आधारभूत ग्रन्थों (प्रस्थानत्रयी) का विस्तृत परिचय दीजिए। \pts{10}

\medskip\rule{\linewidth}{0.2pt}\medskip

\noindent \textbf{Question 2.} अद्वैत वेदान्त में 'अध्यास' की अवधारणा की सोदाहरण व्याख्या कीजिए। \pts{10}

\medskip\rule{\linewidth}{0.2pt}\medskip

\noindent \textbf{Question 3.} श्रीमद्भगवद्गीता के अनुसार निष्काम कर्मयोग का स्वरूप स्पष्ट कीजिए। \pts{10}

\medskip\rule{\linewidth}{0.2pt}\medskip

\noindent \textbf{Question 4.} वेदान्तसार के अनुसार साधन-चतुष्टय का विस्तृत वर्णन कीजिए। \pts{10}

\medskip\rule{\linewidth}{0.2pt}\medskip

\noindent \textbf{Question 5.} अज्ञान के स्वरूप तथा उसकी आवरण एवं विक्षेप शक्तियों की सोदाहरण विवेचना कीजिए। \pts{10}

\bigskip
\rule{\linewidth}{0.2pt}
\bigskip

\begin{center}
  {\large\bfseries SECTION B / खण्ड-ब}\\[2pt]
  {\normalsize (Short Answer Type Questions / लघु उत्तरीय प्रश्न)}
\end{center}

\noindent \textit{सूचना : किन्हीं चार प्रश्नों के उत्तर दीजिए। प्रत्येक प्रश्न 5 अंकों का है।}\\[2pt]

\bigskip

\noindent \textbf{Question 6.} 'सच्चिदानन्द' ब्रह्म के स्वरूप को स्पष्ट कीजिए। \pts{5}

\medskip\rule{\linewidth}{0.2pt}\medskip

\noindent \textbf{Question 7.} जीवनमुक्ति तथा विदेहमुक्ति में अंतर स्पष्ट कीजिए। \pts{5}

\medskip\rule{\linewidth}{0.2pt}\medskip

\noindent \textbf{Question 8.} महावाक्य 'तत्त्वमसि' का वेदान्तीय अर्थ स्पष्ट कीजिए। \pts{5}

\medskip\rule{\linewidth}{0.2pt}\medskip

\noindent \textbf{Question 9.} सूक्ष्म शरीर (पञ्चदश कलाओं) का संक्षिप्त वर्णन कीजिए। \pts{5}

\medskip\rule{\linewidth}{0.2pt}\medskip

\noindent \textbf{Question 10.} विवर्तवाद तथा परिणामवाद में भेद निरूपित कीजिए। \pts{5}

\medskip\rule{\linewidth}{0.2pt}\medskip

\noindent \textbf{Question 11.} पञ्चीकरण प्रक्रिया का संक्षिप्त परिचय दीजिए। \pts{5}

\bigskip
\rule{\linewidth}{0.2pt}
\bigskip

\begin{center}
  {\large\bfseries SECTION C / खण्ड-स}\\[2pt]
  {\normalsize (Very Short Answer Type Questions / अतिलघु उत्तरीय प्रश्न)}
\end{center}

\noindent \textit{सूचना : निम्नलिखित सभी दस प्रश्नों का उत्तर एक-दो वाक्यों में दीजिए। प्रत्येक प्रश्न 2 अंकों का है।}\\[2pt]

\bigskip

\noindent \textbf{Question 12.}
\begin{parts}
  \item प्रस्थानत्रयी के अन्तर्गत कौन-कौन से तीन ग्रन्थ आते हैं? \pts{2}
  \item वेदान्तसार के रचयिता कौन हैं? \pts{2}
  \item ब्रह्मसूत्र के रचयिता का नाम लिखिए। \pts{2}
  \item 'अहं ब्रह्मास्मि' किस उपनिषद् का महावाक्य है? \pts{2}
  \item अज्ञान की कितनी शक्तियाँ मानी गई हैं? \pts{2}
  \item अद्वैत वेदान्त के प्रवर्तक आचार्य कौन हैं? \pts{2}
  \item 'माया' का लक्षण क्या है? \pts{2}
  \item साधन-चतुष्टय के प्रथम साधन का नाम लिखिए। \pts{2}
  \item 'सच्चिदानन्द' में तीन पदों का क्या अर्थ है? \pts{2}
  \item वेदान्त दर्शन का मूल उपजीव्य कौन-सा साहित्य है? \pts{2}
\end{parts}

\vfill
\rule{\linewidth}{0.4pt}

\begin{center}\small
\href{https://bhu-exam.vercel.app/}{Click here to visit our website}\\[4pt]
\href{https://me-aryan.vercel.app/}{Click here to visit our contributor}
\end{center}

\end{document}
